\section{Anwendungen und Zweck}
\label{julia:zweck}
\kopfrechts{Anwendungen und Zweck}%

Die Definitionen für die Julia- und Fatou-Mengen sowie die Mandelbrot-Menge basieren auf interessanten Überlegungen, doch der grundlegende Zweck kann mitunter etwas in den Hintergrund geraten.
Am Anfang des Kapitels wurde mit der Gleichung \eqref{julia:beispielgleichung} und dem Newton-Algorithmus die Überlegung der Funktionen-Iteration motiviert.
Die Iterationsformel \eqref{julia:newton_alg} dieser Gleichung hat zwei Komponenten: die Variable $x_n$ beim Schritt $n$ und den Parameter $a$.
Beide Komponenten lassen sich direkt mit der entwickelten Theorie in Verbindung bringen.

\subsection{Die Variable und die Julia-Mengen}
Für einen fixen Parameterwert $a$ hat die Gleichung \eqref{julia:beispielgleichung} zwei analytische Lösungen: $\!\sqrt{-a}$ und $-\!\sqrt{-a}$.
Weiter ergibt sich aus der Iterationsformel \eqref{julia:newton_alg} die Funktion
\begin{equation*}
    g(x_n) = x_n - \frac{x_n^2 + a}{2x_n} = \frac{x_n^2 - a}{2x_n}.
\end{equation*}
Bei diesen analytischen Lösungen handelt es sich, aus Sicht der Julia-Menge, um Fixpunkte, denn
\begin{equation*}
    g(\!\sqrt{-a}) = \!\sqrt{-a} \qquad \text{und} \qquad g(-\!\sqrt{-a}) = -\!\sqrt{-a}.
\end{equation*}

Möchte man wissen, um welche Art von Fixpunkten es sich handelt, müssen die Werte des Multiplikators \eqref{julia:lambda_multiplicator} für die beiden analytischen Lösungen berechnet werden.
Dazu benötigt man die Ableitung
\begin{equation*}
    \left(g^n\right)^\prime(x_0) = \prod_{k=0}^{n-1} g^\prime\left(g^k(x_0)\right) = g^\prime(x_0) \cdot g^\prime(g(x_0)) \cdot g^\prime(g(g(x_0))) \cdot \ldots
\end{equation*}
und muss diese für die beiden analytischen Lösungen auswerten.
Interessant ist der Faktor $g^\prime(x_0)$, denn sowohl $g^\prime(\!\sqrt{-a})$ als auch $g^\prime(-\!\sqrt{-a})$ sind gleich $0$, wodurch die gesamten Produkte und somit die Multiplikatoren zu $0$ werden.
Gemäss der Definition \ref{julia:lambda_defs} sind daher beide analytischen Lösungen stark anziehende Fixpunkte.

Die komplexe Ebene (oder hier die reellen Zahlen) ist demzufolge in \emph{Einzugsgebiete} unterteilt:
\index{Einzugsgebiet}%
\begin{itemize}
    \item Punkte, für die die Iteration von $g$ zu $\!\sqrt{-a}$ konvergiert
    \item Punkte, für die die Iteration von $g$ zu $-\!\sqrt{-a}$ konvergiert
    \item Punkte, für die die Iteration von $g$ sich anders verhält (z.B. $x = 0$ führt zu Division durch $0$)
\end{itemize}
Die Julia-Menge der Funktion $g$ entspricht genau dem Rand, der diese Einzugsgebiete trennt.
Wählt man einen Punkt auf dem Rand und variiert diesen leicht, fällt man je nach Art der Änderung in unterschiedliche Einzugsgebiete.
Diese Randpunkte verhalten sich also wieder chaotisch.

Möchte man nun die Gleichung \eqref{julia:beispielgleichung} mit der Newton-Iterationsformel \eqref{julia:newton_alg} lösen, benötigt man einen Startwert $x_0$.
Aus Analysis I ist bekannt, dass der Startwert ``genügend nahe'' bei einer Lösung der Gleichung liegen muss.
Berechnet man nun die Julia-Menge der Funktion $g$, findet man so die Einzugsgebiete der Lösungen und kann damit quantifizieren, was ``genügend nahe'' bedeutet.

\subsection{Der Parameter und die Mandelbrot-Menge}
Fixiert man statt dem Parameter, die Variable bei $x_0$, kann so studiert werden, welche Parameterwerte zu welcher Art von Verhalten führen.
Man untersucht nun das Verhalten der Funktion
\begin{equation*}
    g(x_0)_a = \frac{x_0^2 - a}{2x_0}
\end{equation*}
für verschiedene Werte von $a$.
Dies ist analog zur Definition der Mandelbrot-Menge in Abschnitt \ref{julia:sec:mandelbrot} mit dem Unterschied, dass nicht die quadratische Funktion untersucht wird.

Findet man alle Parameterwerte, für die die Rekursion von $g$ beschränkt bleibt, erhält man zwar nicht die klassische Mandelbrot-Menge, aber trotzdem etwas was die gleichen Eigenschaften besitzt.
Hiermit kann also untersucht werden, wie sich verschiedene Parameterwerte $a$ auf das Lösen der Gleichung \eqref{julia:beispielgleichung} mit der Newton-Iterationsformel \eqref{julia:newton_alg} auswirken.
